% future/htm.tex
% mainfile: ../perfbook.tex
% SPDX-License-Identifier: CC-BY-SA-3.0

\section{硬件事务内存}
\label{sec:future:Hardware Transactional Memory}
%
\epigraph{在自动化之前，确保你的报告系统合理地清晰和高效。
	  否则，你的新电脑只会加速这团混乱。}
	 {Robert Townsend}
% If at first you do succeed---try to hide your astonishment.
% Harry F.~Banks

截至2021年，\IXacrf{htm}已经在多种商用通用计算机系统上可用多年~\cite{Yoo:2013:PEI:2503210.2503232,RickMerrit2011PowerTM,ChristianJacobi2012MainframeTM,TimothyHayes2020ARM-HTM}。
本节试图确定HTM在并行程序员工具箱中的位置。

从概念上看，HTM利用处理器缓存和推测执行，
使一组指定的语句（即"事务"）从其他处理器上运行的
任何其他事务的角度来看是原子生效的。
该事务由begin-transaction机器指令启动，
并由commit-transaction机器指令完成。
通常还有一条abort-transaction机器指令，
它会取消推测（如同begin-transaction指令及其后续所有指令都未执行过一样），
并从失败处理程序处开始执行。
失败处理程序的位置通常由begin-transaction指令指定，
可以是显式的失败处理程序地址，也可以通过该指令本身设置的条件码来指定。
每个事务相对于所有其他事务原子执行。

HTM具有多项重要优势，包括自动动态划分数据结构、
减少同步原语的缓存未命中（cache miss），以及支持相当数量的实际应用。

然而，阅读细则总是值得的，HTM也不例外。
本节的一个重点是确定在什么条件下HTM的优势超过其细则中隐藏的复杂性。
为此，\cref{sec:future:HTM Benefits WRT Locking}
描述了HTM的优势，
\cref{sec:future:HTM Weaknesses WRT Locking}描述了其弱点。
这与早期论文~\cite{McKenney2007PLOSTM,PaulEMcKenney2010OSRGrassGreener}
以及前一节中使用的方法相同。\footnote{
	我衷心感谢与其他作者Maged Michael、Josh Triplett
	和Jonathan Walpole以及Andi Kleen进行的许多富有启发的讨论。}

\Cref{sec:future:HTM Weaknesses WRT Locking When Augmented}接着描述了
HTM相对于Linux内核（以及许多用户空间应用程序）中使用的同步原语组合的弱点。
\Cref{sec:future:Where Does HTM Best Fit In?}探讨了HTM
最适合在并行程序员工具箱中的位置，
\cref{sec:future:Potential Game Changers}列出了一些可能
大大增加HTM适用范围和吸引力的事件。
最后，\cref{sec:future:Conclusions}
给出了总结性评述。

\subsection{HTM 相对于锁的优势}
\label{sec:future:HTM Benefits WRT Locking}

HTM的主要优势是：
（1）避免了其他同步原语经常导致的缓存未命中，
（2）能够动态划分数据结构，
以及（3）它具有相当数量的实际应用。
我打破TM的传统，没有单独列出易用性，原因有二。
首先，易用性应该源于HTM的主要优势，这正是本节的重点。
其次，围绕着测试原始编程天赋的尝试存在相当大的争议~\cite{RichardBornat2006SheepGoats,SaeedDehnadi2009SheepGoats,ElizabethPatitsas2020GradesNotBimodal}，
甚至在工作面试中使用小型编程练习也存在争议~\cite{RegBraithwaite2007FizzBuzz}。
这表明我们确实没有牢固地掌握是什么使编程变得容易或困难。
因此，本节的其余部分重点讨论上面列出的三个优势。

\subsubsection{避免同步 Cache Miss}
\label{sec:future:Avoiding Synchronization Cache Misses}

大多数同步机制基于由原子指令操作的数据结构。
因为这些原子指令通常首先需要使相关的\IX{cache line}被运行它们的CPU所拥有，
随后在其他CPU上执行同一同步原语实例将导致缓存未命中。
这些通信缓存未命中严重降低了传统同步机制的性能和可扩展性~\cite[Section 4.2.3]{Anderson97}。

相比之下，HTM通过使用CPU的缓存进行同步，避免了对单独的同步数据结构
及其导致的缓存未命中的需要。
HTM的优势在锁数据结构被放置在单独的缓存行中的情况下最大，
在这种情况下，将给定的临界区转换为HTM事务可以减少一次完整的缓存未命中开销。
这些节省对于短临界区的常见情况可能非常显著，至少在被省略的锁
没有与该锁保护的频繁写入变量共享缓存行的情况下如此。

\QuickQuiz{
	为什么频繁写入的变量与锁变量共享缓存行会有影响？
}\QuickQuizAnswer{
	如果锁与它保护的某些变量在同一缓存行中，
	那么一个CPU对这些变量的写入将使所有其他CPU的该缓存行失效。
	这些\IXpl{invalidation}将产生大量冲突和重试，甚至可能
	降低性能和可扩展性，比使用锁还差。
}\QuickQuizEnd

\subsubsection{数据结构的动态分区}
\label{sec:future:Dynamic Partitioning of Data Structures}

使用某些传统同步机制的一个主要障碍是需要静态划分数据结构。
有许多数据结构可以轻松划分，最突出的例子是哈希表，
其中每个哈希链构成一个分区。
为每个哈希链分配一个锁，就可以轻松地将哈希表中局限于给定链的操作并行化。\footnote{
	同样容易将此方案扩展到访问多个哈希链的操作，
	只需让这些操作按哈希顺序获取所有相关链的锁即可。}
对于数组、基数树、跳表和其他一些数据结构，划分同样很简单。

然而，对于许多类型的树和图，划分相当困难，结果通常也相当复杂~\cite{Ellis80}。
虽然可以使用两阶段锁和哈希锁数组来划分一般数据结构，
但其他技术已被证明更可取~\cite{DavidSMiller2006HashedLocking}，
如将在\cref{sec:future:HTM Weaknesses WRT Locking When Augmented}中讨论的。
鉴于其避免了同步缓存未命中的特性，
HTM因此对于大型不可划分的数据结构是一个非常现实的选择，
至少在更新相对较小的情况下是如此。

\QuickQuiz{
	为什么相对较小的更新对\IXacr{htm}的性能和可扩展性很重要？
}\QuickQuizAnswer{
	更新越大，冲突的概率就越大，因此重试的概率也越大，
	这会降低性能。
}\QuickQuizEnd

\subsubsection{实际价值}
\label{sec:future:Practical Value}

HTM的实际价值已在多个硬件平台上得到了一定程度的证明，包括
Sun Rock~\cite{DaveDice2009ASPLOSRockHTM}、
Azul Vega~\cite{CliffClick2009AzulHTM}、
IBM Blue Gene/Q~\cite{RickMerrit2011PowerTM}、
Intel Haswell TSX~\cite{RaviRajwar2012TSX}和
IBM System z~\cite{ChristianJacobi2012MainframeTM}。

预期的实际优势包括：

\begin{enumerate}
\item	内存数据访问和更新的锁省略~\cite{Martinez01a,Rajwar02a}。
\item	对大型不可划分数据结构的并发访问和小规模随机更新。
\end{enumerate}

然而，HTM也有一些非常实际的不足，将在下一节中讨论。

\subsection{HTM 相对于锁的劣势}
\label{sec:future:HTM Weaknesses WRT Locking}

HTM的概念非常简单：
一组对内存的访问和更新原子地发生。
然而，与许多简单的想法一样，当你将其应用于现实世界中的真实系统时，
就会出现复杂情况。
这些复杂情况如下：

\begin{enumerate}
\item	事务大小限制。
\item	冲突处理。
\item	中止和回滚。
\item	缺乏前进进度保证。
\item	不可撤销操作。
\item	语义差异。
\end{enumerate}

以下各节将分别讨论这些复杂情况，然后给出总结。

\subsubsection{事务大小限制}
\label{sec:future:Transaction-Size Limitations}

当前HTM实现的事务大小限制源于使用处理器缓存来保存
受事务影响的数据。
虽然这允许给定的CPU通过在其缓存范围内执行事务来使事务对其他CPU
显得原子，但这也意味着任何不适合的事务都无法提交。
此外，改变执行上下文的事件，如中断、系统调用、异常、陷阱和上下文切换，
要么必须中止相关CPU上正在进行的事务，要么必须由于其他执行上下文的
缓存占用而进一步限制事务大小。

当然，现代CPU往往具有大型缓存，许多事务所需的数据可以轻松放入
一兆字节的缓存中。
不幸的是，对于缓存来说，单纯的大小并不是唯一重要的因素。
问题在于大多数缓存可以被认为是用硬件实现的哈希表。
然而，硬件缓存不会链接其桶（通常称为\emph{组}），而是
为每组提供固定数量的缓存行。
给定缓存中每组提供的元素数量
称为该缓存的\emph{\IXalt{关联度}{cache associativity}}。

虽然缓存关联度各不相同，但我正在打字的这台笔记本电脑上
L0缓存的八路关联并不罕见。
这意味着如果一个给定的事务需要访问九个缓存行，
且所有九个缓存行都映射到同一组，
那么该事务不可能完成，不管缓存中还有多少兆字节的额外空间。
是的，对于给定数据结构中随机选择的数据元素，
该事务能够提交的概率相当高，
但不能保证~\cite{PaulEMcKenney2012HTMCacheGeometry}。

已有一些研究工作试图缓解这一限制。
全关联\emph{victim缓存}可以缓解关联度约束，
但目前对victim缓存的大小有严格的性能和能效约束。
也就是说，用于未修改缓存行的HTM victim缓存可以相当小，
因为它们只需要保留地址：
数据本身可以写回内存或被其他缓存影子保存，
而地址本身就足以检测冲突写入~\cite{RaviRajwar2012TSX}。

\IXAcrmfst{utm}
方案~\cite{CScottAnanian2006,KevinEMoore2006}
使用DRAM作为极大的victim缓存，但将这些方案集成到
生产级\IXalt{cache-coherence}{cache coherence}机制中
仍然是一个未解决的问题。
此外，使用DRAM作为victim缓存可能会带来不利的
性能和能效后果，特别是如果victim缓存要求
\IXalth{全关联}{fully associative}{cache}的话。
最后，UTM的"无界"方面假设所有DRAM
都可以用作victim缓存，而实际上分配给给定CPU的
大量但仍然固定的DRAM将限制该CPU事务的大小。
其他方案使用硬件和软件transactional memory的组合~\cite{SanjeevKumar2006}，
人们可以想象使用\IXacr{stm}作为HTM的后备机制。

然而，据我所知，除了缩略表示TM读集之外，
当前可用的系统没有实现这些研究想法中的任何一个，
也许是有充分理由的。

\subsubsection{冲突处理}
\label{sec:future:Conflict Handling}

第一个复杂因素是\emph{冲突}的可能性。
例如，假设事务~A和~B定义如下：

\begin{VerbatimU}
Transaction A       Transaction B

x = 1;              y = 2;
y = 3;              x = 4;
\end{VerbatimU}

假设每个事务在各自的处理器上并发执行。
如果事务~A在事务~B存储\co{y}的同时存储\co{x}，
则两个事务都无法继续。
为了理解这一点，假设事务~A执行对\co{y}的存储。
那么事务~A将交错在事务~B内部，违反了事务之间
原子执行的要求。
允许事务~B执行对\co{x}的存储同样违反原子执行要求。
这种情况被称为\emph{冲突}，当两个并发事务访问同一变量
且至少有一个访问是存储时就会发生。
因此系统有义务中止一个或两个事务以允许执行继续。
究竟中止哪个事务是一个有趣的课题，
很可能在未来相当长时间内仍能产出博士论文，
例如参见~\cite{EgeAkpinar2011HTM2TLE}。\footnote{
	Liu和Spear题为"Toxic
	Transactions"~\cite{YujieLiu2011ToxicTransactions}的论文
	特别具有教育意义。}
就本节而言，我们可以假设系统做出随机选择。

另一个复杂因素是冲突检测，至少在最简单的情况下比较直接。
当处理器正在执行事务时，它会标记该事务触及的每个缓存行。
如果处理器的缓存收到涉及已被当前事务标记为已触及的缓存行的请求，
则发生了潜在冲突。
更复杂的系统可能会尝试将当前处理器的事务排序在发送请求的处理器的事务之前，
优化此过程也可能在相当长时间内产出博士论文。
然而本节假设一种非常简单的冲突检测策略。

然而，要使HTM有效工作，冲突概率必须相当低，这反过来要求数据结构的组织方式
能维持足够低的冲突概率。
例如，具有简单插入、删除和搜索操作的红黑树符合此描述，
但维护准确元素计数的红黑树则不符合。\footnote{
	更新计数的需要将导致树的添加和删除操作相互冲突，
	产生强不可交换性~\cite{HagitAttiya2011LawsOfOrder,Attiya:2011:LOE:1925844.1926442,PaulEMcKenney2011SNC}。}
另一个例子是，在单个事务中枚举树中所有元素的红黑树将具有很高的冲突概率，
从而降低性能和可扩展性。
因此，许多串行程序在HTM能有效工作之前需要一些重构。
在某些情况下，实践者会倾向于采取额外步骤
（在红黑树的情况下，也许切换到可划分的数据结构如基数树或哈希表），
直接使用锁，特别是在HTM在所有相关架构上广泛可用之前~\cite{CliffClick2009AzulHTM}。

\QuickQuiz{
	无论选择什么同步机制，红黑树怎么可能高效地枚举树中的所有元素？
}\QuickQuizAnswer{
	在许多情况下，枚举不需要精确。
	在这些情况下，可以使用hazard pointer或\IXacr{rcu}来保护读者，
	这提供了与任何给定插入或删除操作的低冲突概率。
}\QuickQuizEnd

此外，并发事务之间潜在的冲突访问可能导致失败。
处理此类失败将在下一节中讨论。

\subsubsection{中止和回滚}
\label{sec:future:Aborts and Rollbacks}

因为任何事务都可能在任何时候被中止，
所以事务不能包含无法回滚的语句，这一点很重要。
这意味着事务不能执行I/O、系统调用或调试断点
（HTM事务中不能在调试器中单步执行！！！）。
相反，事务必须将自己限制在访问正常的缓存内存。
此外，在某些系统上，中断、异常、陷阱、TLB未命中和其他事件也会中止事务。
鉴于因不当处理错误条件而导致的大量错误，
有理由问中止和回滚对易用性有什么影响。

\QuickQuiz{
	但为什么调试器不能通过在事务的连续行上设置断点来模拟单步执行，
	依赖重试来重新追踪事务早期实例的步骤？
}\QuickQuizAnswer{
	这种方案可能以相当高的概率工作，但它
	可能以令大多数用户相当惊讶的方式失败。
	为了理解这一点，考虑以下事务：

\begin{fcvlabel}[ln:future:htm:debug rollbacks]
\begin{VerbatimN}[commandchars=\\\[\]]
begin_trans();
if (a) {
	do_one_thing();
	do_another_thing();	\lnlbl[another]
} else {
	do_a_third_thing();
	do_a_fourth_thing();
}
end_trans();
\end{VerbatimN}
\end{fcvlabel}

	\begin{fcvref}[ln:future:htm:debug rollbacks]
	假设用户在\clnref{another}处设置了断点，
	该断点触发，中止事务并进入调试器。
	\end{fcvref}
	假设在断点触发和调试器停止所有线程之间，
	某个其他线程将\co{a}的值设为零。
	当可怜的用户试图单步执行程序时，惊讶地发现
	程序现在在else子句中而不是then子句中。

	这\emph{不是}我所说的易于使用的调试器。
}\QuickQuizEnd

当然，中止和回滚引发了HTM是否可用于硬实时系统的问题。
HTM的性能优势是否超过了中止和回滚的代价，如果是，在什么条件下？
事务可以使用优先级提升吗？
还是高优先级线程的事务应该优先中止低优先级线程的事务？
如果是这样，如何有效地通知硬件优先级？
关于HTM实时使用的文献相当稀少，
也许是因为在非实时环境中使HTM良好工作就已经有足够多的问题了。

因为当前的HTM实现可能确定性地中止给定事务，
软件必须提供后备代码。
这种后备代码必须使用某种其他形式的同步，例如锁。
如果曾经使用过基于锁的后备，那么锁的所有限制，
包括\IX{deadlock}的可能性，都会重新出现。
当然可以希望后备不经常使用，这可能允许使用更简单、
更不容易死锁的锁设计。
但这引发了系统如何从使用基于锁的后备过渡回事务的问题。\footnote{
	应用程序卡在后备模式中的可能性被称为"旅鼠效应"，
	这个术语归功于Dave Dice。}
一种方法是使用test-and-test-and-set策略~\cite{Martinez02a}，
让所有人都等到锁释放，允许系统在那个时点以事务模式从干净的状态开始。
然而，这可能导致相当多的自旋，如果锁持有者已被阻塞或被抢占，
这可能不是明智之举。
另一种方法是允许事务与持有锁的线程并行进行~\cite{Martinez02a}，
但这在维护原子性方面带来了困难，
特别是如果线程持有锁的原因是因为相应的事务无法放入缓存。

最后，处理中止和回滚的可能性似乎给开发者增加了额外的负担，
他们必须正确处理所有可能的错误条件组合。

显然，HTM的用户必须在测试后备代码路径和从后备代码
过渡回事务代码方面投入大量验证工作。
也没有理由认为HTM硬件的验证要求不那么艰巨。

\subsubsection{缺乏前进保证}
\label{sec:future:Lack of Forward-Progress Guarantees}

尽管事务大小、冲突和中止/回滚都可能导致事务中止，
人们可能希望足够小且持续时间短的事务可以被保证最终成功。
这将允许无条件重试事务，就像\IXacrmf{cas}和load-linked/store-conditional
（LL/SC）操作在使用这些指令实现原子操作（atomic operation）的代码中被无条件重试一样。

不幸的是，除了低时钟频率的学术研究原型~\cite{MartinSchoeberl2010realtimeTM}之外，
当前可用的HTM实现拒绝做出任何形式的前进进度保证。
如前所述，HTM因此不能用于在这些系统上避免死锁。
希望未来的HTM实现将提供某种形式的前进进度保证。
在此之前，HTM在实时应用中必须极其谨慎地使用。

截至2021年，这一暗淡图景的唯一例外是IBM大型机，
它提供了\emph{受约束事务}~\cite{ChristianJacobi2012MainframeTM}。
约束条件相当严格，在\cref{sec:future:Forward-Progress Guarantees}中介绍。
看看HTM前进进度保证是否会从大型机迁移到通用CPU系列将是有趣的。

\subsubsection{不可撤销操作}
\label{sec:future:Irrevocable Operations}

中止和回滚的另一个后果是HTM事务无法容纳不可撤销的操作。
当前的HTM实现通常通过要求事务中的所有访问都是对可缓存内存的访问
（从而禁止MMIO访问）并在中断、陷阱和异常时中止事务
（从而禁止系统调用）来执行此限制。

请注意，只要缓冲区填充/刷新操作在事务外发生，
HTM事务可以容纳缓冲I/O。
这之所以有效，是因为向缓冲区添加数据和从中移除数据是可撤销的：
只有实际的缓冲区填充/刷新操作是不可撤销的。
当然，这种缓冲I/O方法的效果是将I/O包含在事务的占用空间中，
增加了事务的大小，从而增加了失败的概率。

\subsubsection{语义差异}
\label{sec:future:Semantic Differences}

虽然HTM在许多情况下可以作为锁的直接替代品使用
（因此得名\IXacrfst{tle}~\cite{DaveDice2008TransactLockElision}），
但语义上存在细微差异。
Blundell~\cite{Blundell2006TMdeadlock}给出了一个特别棘手的例子，
涉及协调的基于锁的临界区，在事务方式执行时会导致死锁或\IX{livelock}，
但一个更简单的例子是空临界区。

在基于锁的程序中，一个空临界区将保证
之前持有该锁的所有进程现在都已释放它。
2.4版Linux内核的网络栈使用了这种习惯用法来协调配置更改。
但如果将这个空临界区转换为事务，结果就是一个空操作。
所有先前临界区已终止的保证将丢失。
换句话说，事务锁省略保留了锁的数据保护语义，
但丢失了锁的基于时间的消息传递语义。

\QuickQuizSeries{%
\QuickQuizB{
	但为什么\emph{任何人}会需要一个空的基于锁的临界区？
}\QuickQuizAnswerB{
	参见\cref{sec:locking:Exclusive Locks}中
	\QuickQuizARef{\QlockingQemptycriticalsection}的答案。

	然而，有人声称，给定一个没有前进进度保证的强原子\IXacr{htm}
	实现，任何基于空临界区的基于内存的锁设计
	在事务锁省略存在的情况下将正确运行。
	虽然我没有看到这个声明的证明，但有一个直接的理由。
	主要思想是，在强原子HTM实现中，
	给定事务的结果在事务成功完成之前是不可见的。
	因此，如果你能看到一个事务已经开始，
	则保证它已经完成，这意味着
	后续的空锁临界区将成功地"等待"它——毕竟，
	不需要等待。

	这种推理不适用于弱原子系统
	（包括许多\IXacr{stm}实现），也不适用于
	使用内存以外的方式进行通信的基于锁的程序。
	这样的方式之一是时间的流逝（例如，在硬实时系统中）
	或优先级的流动（例如，在软实时系统中）。

	依赖优先级提升的锁设计特别值得关注。
}\QuickQuizEndB
%
\QuickQuizM{
	事务锁省略难道不能通过简单地选择不省略空的基于锁的临界区
	来轻松处理锁的基于时间的消息传递语义吗？
}\QuickQuizAnswerM{
	它可以这样做，但这既不必要也不充分。

	在空临界区由条件编译导致的情况下，这是不必要的。
	在这种情况下，锁的唯一目的很可能是保护数据，
	因此完全省略它是正确的做法。
	实际上，保留空的基于锁的临界区会降低性能和可扩展性。

	另一方面，非空的基于锁的临界区可能同时依赖于
	锁的数据保护和基于时间的消息传递语义。
	在这种情况下使用事务锁省略将是不正确的，
	并会导致错误。
}\QuickQuizEndM
%
\QuickQuizE{
	鉴于现代硬件~\cite{PeterOkech2009InherentRandomness}，
	怎么能指望依赖时序的并行软件正常工作呢？
}\QuickQuizAnswerE{
	简短的回答是，在普通的商用硬件上，
	基于任何精细粒度时序的同步设计都是不明智的，
	不能指望在所有条件下正确运行。

	话虽如此，有些专为硬实时使用而设计的系统要确定性得多。
	在你使用这样的系统的（非常不可能的）情况下，
	这里有一个展示基于时间的同步如何工作的简单例子。
	再次强调，\emph{不要}在商用微处理器上尝试这种方法，
	因为它们具有高度不确定性的性能特征。

	这个例子使用多个工作线程和一个控制线程。
	每个工作线程对应一个出站数据流，
	在每个工作单元执行后将当前时间
	（例如，从\co{clock_gettime()}系统调用获取）
	记录在每线程的\co{my_timestamp}变量中。
	此示例的实时性质导致了以下约束集：

	\begin{enumerate}
	\item	如果给定的工作线程未能在超过
		\co{MAX_LOOP_TIME}的时间段内更新其时间戳，
		则为致命错误。
	\item	锁被谨慎地用于访问和更新全局状态。
	\item	锁在给定线程优先级内按严格的FIFO顺序授予。
	\end{enumerate}

	当工作线程完成其数据流时，它们必须与应用程序的其余部分解除关联，
	并在初始化为\co{-1}的每线程\co{my_status}变量中
	放置一个状态值。
	线程不退出；它们被放在线程池中以适应后续的处理需求。
	控制线程根据需要分配（和重新分配）工作线程，
	并维护线程状态的直方图。
	控制线程以不高于工作线程的实时优先级运行。

	工作线程的代码如下：

\begin{VerbatimN}
	int my_status = -1;  /* Thread local. */

	while (continue_working()) {
		enqueue_any_new_work();
		wp = dequeue_work();
		do_work(wp);
		my_timestamp = clock_gettime(...);
	}

	acquire_lock(&departing_thread_lock);

	/*
	 * Disentangle from application, might
	 * acquire other locks, can take much longer
	 * than MAX_LOOP_TIME, especially if many
	 * threads exit concurrently.
	 */
	my_status = get_return_status();
	release_lock(&departing_thread_lock);

	/* thread awaits repurposing. */
\end{VerbatimN}

	控制线程的代码如下：

\begin{fcvlabel}[ln:future:htm:control thread]
\begin{VerbatimN}[commandchars=\\\@\$]
	for (;;) {
		for_each_thread(t) {
			ct = clock_gettime(...);
			d = ct - per_thread(my_timestamp, t);
			if (d >= MAX_LOOP_TIME) {	\lnlbl@if$
				/* thread departing. */	\lnlbl@dep:b$
				acquire_lock(&departing_thread_lock); \lnlbl@acq$
				release_lock(&departing_thread_lock); \lnlbl@rel$
				i = per_thread(my_status, t);
				status_hist[i]++; /* Bug if TLE! */ \lnlbl@dep:e$
			}
		}
		/* Repurpose threads as needed. */
	}
\end{VerbatimN}
\end{fcvlabel}

	\begin{fcvref}[ln:future:htm:control thread]
	\Clnref{if}使用时间的流逝来推断线程已退出，
	如果是这样则执行\clnref{dep:b,dep:e}。
	\clnref{acq,rel}上的空锁临界区保证任何正在退出过程中的线程
	已经完成（记住锁是按FIFO顺序授予的！）。
	\end{fcvref}

	再次强调，不要在商用微处理器上尝试这种方法。
	毕竟，即使在专门为硬实时使用而设计的系统上，
	正确实现这一点也够困难的了！
}\QuickQuizEndE
}

锁和事务之间的一个重要语义差异是在基于锁的实时程序中
用于避免优先级反转的优先级提升。
优先级反转可能发生的一种方式是当持有锁的低优先级线程
被中等优先级的CPU密集型线程抢占时。
如果每个CPU至少有一个这样的中等优先级线程，
低优先级线程将永远没有机会运行。
如果高优先级线程现在试图获取该锁，它将被阻塞。
它无法获取锁直到低优先级线程释放它，
低优先级线程无法释放锁直到它有机会运行，
而它无法有机会运行直到某个中等优先级线程放弃其CPU。
因此，中等优先级线程实际上阻塞了高优先级进程，
这就是"优先级反转"这个名称的由来。

\begin{listing}
\begin{fcvlabel}[ln:future:Exploiting Priority Boosting]
\begin{VerbatimL}[commandchars=\\\@\$]
void boostee(void)		\lnlbl@low:b$
{
	int i = 0;

	acquire_lock(&boost_lock[i]);	\lnlbl@1stacq$
	for (;;) {
		acquire_lock(&boost_lock[!i]);
		release_lock(&boost_lock[i]);
		i = i ^ 1;
		do_something();
	}
}				\lnlbl@low:e$

void booster(void)		\lnlbl@high:b$
{
	int i = 0;

	for (;;) {
		usleep(500); /* sleep 0.5 ms. */
		acquire_lock(&boost_lock[i]);	\lnlbl@acq$
		release_lock(&boost_lock[i]);	\lnlbl@rel$
		i = i ^ 1;
	}
}                               \lnlbl@high:e$
\end{VerbatimL}
\end{fcvlabel}
\caption{利用优先级提升}
\label{lst:future:Exploiting Priority Boosting}
\end{listing}

避免优先级反转的一种方法是\emph{优先级继承}，
其中被锁阻塞的高优先级线程临时将其优先级捐赠给锁的持有者，
这也被称为\emph{优先级提升}。
然而，优先级提升可以用于避免优先级反转以外的目的，
如\cref{lst:future:Exploiting Priority Boosting}所示。
\begin{fcvref}[ln:future:Exploiting Priority Boosting]
该清单的\clnrefrange{low:b}{low:e}展示了一个必须大约每毫秒运行一次的低优先级进程，
而同一清单的\clnrefrange{high:b}{high:e}展示了一个使用优先级提升
来确保\co{boostee()}按需定期运行的高优先级进程。

\co{boostee()}函数通过始终持有两个\co{boost_lock[]}锁中的一个
来实现这一点，以便\co{booster()}的\clnrefrange{acq}{rel}
可以根据需要提升优先级。
\end{fcvref}

\QuickQuiz{
	但\cref{lst:future:Exploiting Priority Boosting}中的
	\co{boostee()}函数交替以相反的顺序获取锁！
	这不会导致死锁吗？
}\QuickQuizAnswer{
	不会导致死锁。
	要产生死锁，两个不同的线程必须各自以相反的顺序
	获取两个锁，这在此示例中不会发生。
	然而，像lockdep~\cite{JonathanCorbet2006lockdep}
	这样的死锁检测器会将此标记为误报。
}\QuickQuizEnd

\begin{fcvref}[ln:future:Exploiting Priority Boosting]
这种安排要求\co{boostee()}在系统变忙之前在\clnref{1stacq}处获取其第一个锁，
但即使在现代硬件上也很容易安排。

不幸的是，这种安排在事务锁省略存在的情况下可能会失效。
\co{boostee()}函数重叠的临界区变成一个无限事务，
迟早会中止，例如在运行\co{boostee()}函数的线程
第一次被抢占时。
此时\co{boostee()}将回退到使用锁，但鉴于其低优先级
以及安静的初始化期间已经结束
（这毕竟是\co{boostee()}被抢占的原因），
这个线程可能永远不会再有机会运行。

而如果\co{boostee()}线程没有持有锁，那么
\cref{lst:future:Exploiting Priority Boosting}中
\co{booster()}线程在\clnref{acq,rel}上的空临界区
将变成一个没有效果的空事务，
从而\co{boostee()}永远不会运行。
此例说明了transactional memory的回滚重试语义的一些微妙后果。
\end{fcvref}

鉴于经验可能会揭示更多微妙的语义差异，
基于HTM的锁省略在大型程序中的应用应当谨慎进行。
话虽如此，在适用的地方，基于HTM的锁省略可以消除
与锁变量相关的缓存未命中，这在截至2015年初的
大型实际软件系统中已经带来了数十个百分点的性能提升。
因此我们可以期望在提供可靠支持的硬件上看到这种技术的大量使用。

\QuickQuiz{
	所以一群人着手取代锁，结果大多只是在优化锁？
}\QuickQuizAnswer{
	至少他们完成了一些有用的事情！
	也许随着时间的推移，\IXacr{htm}会继续取得进步~\cite{Siakavaras2017CombiningHA,DimitriosSiakavaras2020RCU-HTM-B+Trees,ChristinaGiannoula2018HTM-RCU-graphcoloring,SeongJaePark2020HTMRCUlock}。
}\QuickQuizEnd

\subsubsection{总结}
\label{sec:future:HTM Weaknesses WRT Locking: Summary}

% future/HTMtable.tex
% SPDX-License-Identifier: CC-BY-SA-3.0

\begin{table*}
\centering
% \scriptsize
\small
\input{future/HTMtableColor}
\setlength{\tabcolsep}{4pt}\OneColumnHSpace{-.9in}
\ebresizewidth{
\begin{tabularx}{6.5in}{p{0.95in}cXcX}
\toprule
  &
    & \multicolumn{1}{c}{锁} &
      & \multicolumn{1}{c}{Hardware Transactional Memory} \\
\midrule
  基本思想 &
    & 一次只允许一个线程访问给定的一组对象。 &
      & 使对给定对象集的操作以原子方式执行。 \\
\midrule
  适用范围 &
    & \Pl 处理所有操作。 &
      & \Pl 处理可撤销操作。 \\
\addlinespace[4pt]
  &
    & &
      & \Mn 不可撤销操作迫使回退（通常回退到锁）。 \\
\midrule
  可组合性 &
    & \Dw 受死锁限制。 &
      & \Dw 受不可撤销操作、事务大小和死锁限制
        （假设基于锁的回退代码）。 \\
\midrule
  可扩展性与性能 &
    & \Mn 数据必须可分区以避免锁竞争。 &
      & \Mn 数据必须可分区以避免冲突。 \\
\cmidrule{3-5}
  &
    & \Dw 分区通常必须在设计时固定。 &
      & \Pl 自动动态调整分区，精确到缓存行边界。 \\
\addlinespace[4pt]
  &
    & &
      & \Mn 回退需要分区（对于罕见的回退不太重要）。 \\
\cmidrule{3-5}
  &
    & \Dw 锁原语通常导致昂贵的缓存未命中
      和内存屏障指令。 &
      & \Mn 事务开始/结束指令通常不会导致缓存未命中，
        但确实有内存排序和开销影响。 \\
\cmidrule{3-5}
  &
    & \Pl 竞争效应集中在获取和释放上，因此
      临界区以全速运行。 &
      & \Mn 竞争会中止冲突的事务，即使它们已经
        运行了很长时间。 \\
\cmidrule{3-5}
  &
    & \Pl 私有化操作简单、直观、高效且可扩展。 &
      & \Mn 私有化数据增加事务大小。 \\
\midrule
  硬件支持 &
    & \Pl 通用硬件即可。 &
      & \Mn 需要新硬件（正开始变得可用）。 \\
\cmidrule{3-5}
  &
    & \Pl 性能不受缓存几何细节影响。 &
      & \Mn 性能严重依赖缓存几何结构。 \\
\midrule
  软件支持 &
    & \Pl API已存在，有大量代码和经验，调试器正常工作。 &
      & \Mn API正在出现，DBMS之外经验很少，事务中的
        断点可能有问题。 \\
\midrule
  与其他机制的交互 &
    & \Pl 有长期成功交互的经验。 &
      & \Dw 刚刚开始研究交互。 \\
\midrule
  实际应用 &
    & \Pl 有。 &
      & \Pl 有。 \\
\midrule
  广泛适用性 &
    & \Pl 有。 &
      & \Mn 尚无定论。 \\
\bottomrule
\end{tabularx}
}
\IfTwoColumn{
\caption{锁与HTM的比较（\colorbox{plus}{优势}、
  \colorbox{minus}{劣势}、\colorbox{down}{严重劣势}）}
}{
\caption{锁与HTM的比较（\colorbox{plus}{优势}、
  \colorbox{minus}{劣势}、
  \colorbox{down}{严重} \colorbox{down}{劣势}）}
}
\label{tab:future:Comparison of Locking and HTM}
\end{table*}


虽然HTM似乎可能会有引人注目的用例，
但当前的实现存在严重的事务大小限制、冲突处理复杂性、
中止和回滚问题以及需要仔细处理的语义差异。
HTM相对于锁的当前状况总结在
\cref{tab:future:Comparison of Locking and HTM}中。
可以看出，虽然HTM的当前状态缓解了锁的一些严重不足，\footnote{
	公平地说，重要的是要强调锁的不足确实有广为人知且广泛使用的工程解决方案，
	包括死锁检测器~\cite{JonathanCorbet2006lockdep}、
	大量已适应锁的数据结构，以及如
	\cref{sec:future:HTM Weaknesses WRT Locking When Augmented}中讨论的
	长期增强历史。
	此外，如果锁真的像快速浏览许多学术论文可能合理地使人相信的那样糟糕，
	那些大型的基于锁的并行程序（包括自由开源软件和专有软件）
	到底是从哪里来的呢？}
但它通过引入大量自身的不足来实现这一点。
这些不足被TM社区的领导者所承认~\cite{AlexanderMatveev2012PessimisticTM}。\footnote{
	此外，在2011年初，我受邀对transactional memory背后的
	一些假设进行批评~\cite{PaulEMcKenney2011Verico}。
	听众出奇地不敌视，不过也许他们对我手下留情是因为
	我在做报告时严重倒时差。}

此外，这还不是完整的故事。
锁通常不是单独使用的，而是通常由其他同步机制增强，
包括引用计数、原子操作、非阻塞数据结构、
\IXpl{hazard pointer}~\cite{MagedMichael04a,HerlihyLM02}
和RCU~\cite{McKenney98,McKenney01a,ThomasEHart2007a,PaulEMcKenney2012ELCbattery}。
下一节将探讨这种增强如何改变等式。

\subsection{增强后的 HTM 相对于锁的劣势}
\label{sec:future:HTM Weaknesses WRT Locking When Augmented}

% future/HTMtableRCU.tex
% SPDX-License-Identifier: CC-BY-SA-3.0

\begin{table*}
\centering
\small
\input{future/HTMtableColor}
\setstretch{0.95}
\setlength{\tabcolsep}{4pt}\OneColumnHSpace{-.9in}
\ebresizewidth{
\resizebox{6.5in}{!}{
\begin{tabularx}{6.8in}{p{0.95in}cXcX}
\toprule
  &
    & \multicolumn{1}{c}{用户空间RCU或Hazard Pointer增强的锁} &
      & \multicolumn{1}{c}{Hardware Transactional Memory} \\
\midrule
  基本思想 &
    & 一次只允许一个线程访问给定的一组对象。 &
      & 使对给定对象集的操作以原子方式执行。 \\
\midrule
  适用范围 &
    & \Pl 处理所有操作。 &
      & \Pl 处理可撤销操作。 \\
\addlinespace[4pt]
  &
    & &
      & \Mn 不可撤销操作迫使回退（通常回退到锁）。 \\
\midrule
  可组合性 &
    & \Pl 读取端仅受宽限期等待操作限制。 &
      & \Dw 受不可撤销操作、事务大小和死锁限制
        （假设基于锁的回退代码）。\\
\addlinespace[4pt]
  &
    & \Mn 更新端受死锁限制。读取端减少死锁。 &
      & \\
\midrule
  可扩展性与性能 &
    & \Mn 数据必须可分区以避免更新端之间的锁竞争。 &
      & \Mn 数据必须可分区以避免冲突。 \\
\addlinespace[4pt]
  &
    & \Pl 读取端不需要分区。 &
      & \\
\cmidrule{3-5}
  &
    & \Dw 更新端的分区通常必须在设计时固定。 &
      & \Pl 自动动态调整分区，精确到缓存行边界。 \\
\cmidrule{3-5}
  &
    & \Pl 读取端不需要分区。 &
      & \Mn 回退需要分区（对于罕见的回退不太重要）。 \\
\cmidrule{3-5}
  &
    & \Dw 更新端锁原语通常导致昂贵的缓存未命中
      和内存屏障指令。 &
      & \Mn 事务开始/结束指令通常不会导致缓存未命中，
        但确实有内存排序和开销影响。 \\
\cmidrule{3-5}
  &
    & \Pl 更新端竞争效应集中在获取和释放上，
      因此临界区以全速运行。 &
      & \Mn 竞争会中止冲突的事务，即使它们已经
        运行了很长时间。 \\
\addlinespace[4pt]
  &
    & \Pl 读取端不与更新端或彼此竞争。 &
      & \\
\cmidrule{3-5}
  &
    & \Pl 读取端原语通常是有界无等待的，开销低。
      （Hazard pointer为无锁，开销低。） &
      & \Mn 只读事务受冲突和回滚影响。
        除回退代码提供的之外没有前向进展保证。 \\
\cmidrule{3-5}
  &
    & \Pl 当数据仅对更新端可见时，私有化操作
      简单、直观、高效且可扩展。 &
      & \Mn 私有化数据增加事务大小。 \\
\addlinespace[4pt]
  &
    & \Mn 对于读取端可见的数据，私有化操作开销较大
      （但仍然直观且可扩展）。 &
      & \\
\midrule
  硬件支持 &
    & \Pl 通用硬件即可。 &
      & \Mn 需要新硬件（正开始变得可用）。 \\
\cmidrule{3-5}
  &
    & \Pl 性能不受缓存几何细节影响。 &
      & \Mn 性能严重依赖缓存几何结构。 \\
\midrule
  软件支持 &
    & \Pl API已存在，有大量代码和经验，调试器正常工作。 &
      & \Mn API正在出现，DBMS之外经验很少，事务中的
        断点可能有问题。 \\
\midrule
  与其他机制的交互 &
    & \Pl 有长期成功交互的经验。 &
      & \Dw 刚刚开始研究交互。 \\
\midrule
  实际应用 &
    & \Pl 有。 &
      & \Pl 有。 \\
\midrule
  广泛适用性 &
    & \Pl 有。 &
      & \Mn 尚无定论。 \\
\bottomrule
\end{tabularx}
}}
\caption{锁（由RCU或Hazard Pointer增强）与HTM的比较
  （\colorbox{plus}{优势}、\colorbox{minus}{劣势}、
  \colorbox{down}{严重劣势}）}
\label{tab:future:Comparison of Locking (Augmented by RCU or Hazard Pointers) and HTM}
\end{table*}


实践者长期以来使用引用计数、原子操作、非阻塞数据结构、
hazard pointer和RCU来避免锁的一些不足。
例如，在许多情况下，可以通过使用引用计数、hazard pointer或RCU
来保护数据结构来避免死锁，
特别是对于只读临界区~\cite{MagedMichael04a,HerlihyLM02,MathieuDesnoyers2012URCU,DinakarGuniguntala2008IBMSysJ,ThomasEHart2007a}。
这些方法还减少了划分数据结构的需要，
正如在\cref{chp:Data Structures}中所见。
RCU进一步提供了无竞争的有界无等待读端原语~\cite{McKenney98,MathieuDesnoyers2012URCU}，
而hazard pointer提供了无锁的读端原语~\cite{Michael02a,HerlihyLM02,MagedMichael04a}。
将这些考虑因素添加到\cref{tab:future:Comparison of Locking and HTM}中，
得到了\cref{tab:future:Comparison of Locking (Augmented by RCU or Hazard Pointers) and HTM}
中所示的增强锁与HTM之间的更新比较。
两个表之间差异的总结如下：

\begin{enumerate}
\item	使用非阻塞读端机制缓解了死锁问题。
\item	hazard pointer和RCU等读端机制可以在不可划分的数据上高效操作。
\item	hazard pointer和RCU不会彼此竞争，也不会与更新者竞争，
	允许读多写少工作负载获得出色的性能和可扩展性。
\item	hazard pointer和RCU提供前进进度保证
	（分别为无锁和有界无等待）。
\item	hazard pointer和RCU的私有化操作是直接的。
\end{enumerate}

\IfEbookSize{}{
\input{future/HTMtableFull}

对于视力好的人，
\cref{tab:future:Comparison of Locking (Plain and Augmented) and HTM}
合并了\cref{tab:future:Comparison of Locking and HTM,%
tab:future:Comparison of Locking (Augmented by RCU or Hazard Pointers) and HTM}。
}

\QuickQuiz{
	\Cref{tab:future:Comparison of Locking and HTM,tab:future:Comparison of Locking (Augmented by RCU or Hazard Pointers) and HTM}
	声称硬件才刚刚开始可用。
	但HTM硬件支持不是已经广泛可用近十年了吗？
}\QuickQuizAnswer{
	是也不是。
	看来即使是在真实硬件中实现TM的HTM子集也比看起来要复杂一些~\cite{ChristianJacobi2012MainframeTM,ScottWasson2014HaswellDisableTSX,Intel2020HaswellTSXordering,Intel2021perfTSXordering,MichaelLarabel2021DisableTSX}。
	因此，令人遗憾的事实是，截至2021年"刚刚开始可用"
	这一说法过于准确了。
	实际上，供应商已经开始弃用其HTM实现~\cite[Book III Appendix A]{PowerISA3.1-2020}。
}\QuickQuizEnd

当然，也可以增强HTM，如下一节中讨论的那样。

\subsection{HTM 最适合用在哪里？}
\label{sec:future:Where Does HTM Best Fit In?}

虽然HTM的适用范围可能还需要一段时间才能像
\cpageref{fig:defer:RCU Areas of Applicability}上
\cref{fig:defer:RCU Areas of Applicability}中所示的RCU那样清晰地划定，
但这不是不开始朝这个方向前进的理由。

HTM似乎最适合涉及对在大型多处理器上运行的相对大型内存数据结构的
不同部分进行相对小的更改的更新密集型工作负载，
因为这满足了当前HTM实现的大小限制，同时最小化了冲突
及伴随的中止和回滚的概率。
这种场景也是给定当前同步原语相对难以处理的场景。

将锁与HTM结合使用似乎可以克服HTM在不可撤销操作方面的困难，
而使用RCU或hazard pointer可能会缓解HTM对遍历数据结构大部分内容的
只读操作的事务大小限制~\cite{SeongJaePark2020HTMRCUlock}。
当前的HTM实现会无条件中止与RCU或hazard pointer读者冲突的更新事务，
但也许未来的HTM实现将与这些同步机制更顺畅地互操作。
与此同时，更新与大型RCU或hazard pointer读端临界区冲突的概率
应该远小于与等效只读事务冲突的概率。\footnote{
	相当讽刺的是，严格的事务机制出现在共享内存系统中的时间
	恰好是NoSQL数据库放松传统数据库应用对严格事务依赖的时候。
	尽管如此，HTM确实实现了TM的易用性承诺，
	尽管是用于对Linux内核地址空间随机化防御机制的
	黑帽攻击~\cite{YeongjinJang2016TSXbreakKASLR,Jang:2016:BKA:2976749.2978321}。}
尽管如此，持续不断的RCU或hazard pointer读者完全有可能
因为相应的持续冲突流而使更新者饥饿。
可以通过（以显著的硬件成本和复杂性为代价）
给事务外读取提供被加载内存位置的事务前副本来消除这一漏洞。

HTM事务必须有后备方案的事实可能在某些情况下
会将数据结构的静态可划分性重新强加给HTM。
如果未来的HTM实现提供前进进度保证，
这一限制可能会得到缓解，这可能在某些情况下消除对后备代码的需要，
进而可能允许HTM在冲突概率较高的情况下高效使用。

简而言之，虽然HTM可能有重要的用途和应用，
但它是并行程序员工具箱中的另一个工具，而不是工具箱整体的替代品。

\subsection{潜在的变革因素}
\label{sec:future:Potential Game Changers}

可能大大增加HTM需求的变革因素包括以下几点：

\begin{enumerate}
\item	前进进度保证。
\item	事务大小增加。
\item	改进的调试支持。
\item	弱原子性。
\end{enumerate}

以下各节对此进行了展开讨论。

\subsubsection{前进保证}
\label{sec:future:Forward-Progress Guarantees}

正如\cref{sec:future:Lack of Forward-Progress Guarantees}中所讨论的，
当前的HTM实现缺乏前进保证，这要求必须有回退软件来处理HTM失败。
当然，提出保证要求很容易，但提供保证并不总是那么容易。
对于HTM而言，实现保证的障碍包括缓存大小和关联度、TLB大小和关联度、
事务持续时间和中断频率，以及调度器的实现。

缓存大小和关联度已在\cref{sec:future:Transaction-Size Limitations}中讨论过，
同时还介绍了一些旨在解决当前限制的研究工作。
然而，HTM的前进保证将伴随着大小限制，尽管这些限制可能有朝一日会变得很大。
那么，为什么当前的HTM实现不对小事务提供前进保证呢？例如，将事务限制在
缓存的关联度范围内？
一个可能的原因是需要处理硬件故障。
例如，一个出故障的缓存SRAM单元可能通过停用该故障单元来处理，
从而降低缓存的关联度，因此也降低了能够保证前进的事务的最大尺寸。
鉴于这只是减少了保证的事务大小，看来可能还有其他原因在起作用。
也许在生产级硬件上提供前进保证比人们想象的更加困难，
考虑到在软件中做出前进保证的难度，这是一个完全合理的解释。
将问题从软件转移到硬件并不一定使其更容易解决~\cite{ChristianJacobi2012MainframeTM}。

对于使用物理标记和物理索引的缓存，事务能够放入缓存是不够的。
它的地址转换还必须能放入TLB\@。
因此，任何前进保证还必须考虑TLB的大小和关联度。

鉴于在当前的HTM实现中，中断、陷阱和异常都会中止事务，
因此给定事务的执行持续时间必须短于预期的中断间隔。
无论给定事务接触的数据多么少，如果它运行时间过长，就会被中止。
因此，任何前进保证不仅必须以事务大小为条件，还必须以事务持续时间为条件。

前进保证在很大程度上取决于确定多个冲突事务中哪个应该被中止的能力。
很容易想象一个无尽的事务系列，每个事务中止一个更早的事务，
而自己又被一个更晚的事务中止，以至于没有任何事务能真正提交。
冲突处理的复杂性从已经提出的大量HTM冲突解决策略中可见一斑~\cite{EgeAkpinar2011HTM2TLE,YujieLiu2011ToxicTransactions}。
事务外的访问引入了额外的复杂性，
正如Blundell~\cite{Blundell2006TMdeadlock}所指出的那样。
很容易将所有这些问题归咎于事务外的访问，
但这种思路的荒谬之处可以通过将每个事务外的访问放入其自身的
单次访问事务中来轻松证明。
问题在于访问模式本身，而不是它们是否恰好被封装在事务中。

最后，事务的任何前进保证还取决于调度器，
它必须让执行事务的thread运行足够长的时间来成功提交。

因此，HTM厂商提供前进保证存在重大障碍。
然而，如果其中任何一家做到了，其影响将是巨大的。
这意味着HTM事务将不再需要软件回退，
也就意味着HTM终于可以兑现TM消除死锁的承诺。

然而，在2012年末，IBM大型机宣布了一种HTM实现，
除了通常的尽力而为HTM实现外，还包括\emph{受约束事务}~\cite{ChristianJacobi2012MainframeTM}。
受约束事务以\co{tbeginc}指令开始，
而不是用于尽力而为事务的\co{tbegin}指令。
受约束事务保证总是会（最终）完成，
所以如果事务中止，硬件不是跳转到回退路径
（如同尽力而为事务那样），而是在\co{tbeginc}指令处重新启动事务。

大型机架构师需要采取极端措施来兑现这个前进保证。
如果给定的受约束事务反复失败，CPU可能会禁用分支预测、
强制顺序执行，甚至禁用流水线。
如果反复失败是由于高竞争造成的，CPU可能会禁用推测性取指、
引入随机延迟，甚至序列化冲突CPU的执行。
"有趣的"前进场景涉及少则两个CPU，多则一百个CPU。
也许这些极端措施提供了一些线索，解释了为什么其他CPU
至今没有提供受约束事务。

顾名思义，受约束事务实际上受到严格的约束：

\begin{enumerate}
\item	最大数据占用量是四个内存块，
	其中每个块不能大于32字节。
\item	最大代码占用量是256字节。
\item	如果给定的4K页面包含受约束事务的代码，
	则该页面不能包含该事务的数据。
\item	可执行的最大汇编指令数是32条。
\item	禁止向后跳转。
\end{enumerate}

尽管如此，这些约束支持许多重要的数据结构，
包括链表、栈、队列和数组。
因此，受约束HTM似乎有可能成为并行程序员工具箱中的重要工具。

请注意，这些前进保证不必是绝对的。
例如，假设一种HTM的使用以全局锁作为回退。
假设回退机制已经被仔细设计以避免
\cref{sec:future:Aborts and Rollbacks}中讨论的"旅鼠效应"，
那么如果HTM回滚足够不频繁，全局锁就不会成为瓶颈。
话虽如此，系统越大、临界区越长、
从"旅鼠效应"中恢复所需的时间越长，
"足够不频繁"就需要越罕见。

\subsubsection{事务大小增加}
\label{sec:future:Transaction-Size Increases}

前进保证很重要，但正如我们所看到的，它们将是基于事务大小和持续时间的
有条件保证。
已经取得了一些进展，例如，一些商用HTM实现使用近似技术来支持
极其庞大的HTM读集~\cite{RaviRajwar2012TSX}。
再例如，\Power{8} HTM支持挂起事务，
这可以避免将不相关的访问添加到被挂起事务的读集和
写集中~\cite{Le:2015:TMS:3266491.3266500}。
这一能力已被用于产生高性能的
读写锁~\cite{PascalFelber2016rwlockElision}。

需要注意的是，即使是小尺寸的保证也将非常有用。
例如，两条cache line的保证就足以实现栈、队列或双端队列。
然而，更大的数据结构需要更大的保证，例如，
按顺序遍历一棵树需要与树中节点数量相等的保证。
因此，即使是适度增加保证的大小也会增加HTM的实用性，
从而增加CPU提供这种保证或提供足够好的替代方案的需求。

\subsubsection{改进的调试支持}
\label{sec:future:Improved Debugging Support}

事务大小的另一个抑制因素是调试事务的需求。
当前机制的问题在于，单步异常会中止外层的事务。
对这个问题有许多变通方法，包括模拟处理器（慢！）、
用STM替代HTM（慢而且语义略有不同！）、
使用重复重试来模拟前进的回放技术（奇怪的失败模式！），以及
完全支持调试HTM事务（复杂！）。

如果某个HTM厂商能够生产出允许在事务内直接使用经典调试技术的
HTM系统，包括断点、单步执行和打印语句，
这将使HTM更具吸引力。
一些事务内存（transactional memory）研究者在2013年开始认识到这个问题，
至少有一个提案涉及硬件辅助调试设施~\cite{JustinGottschlich2013TMdebug}。
当然，这个提案依赖于现成的硬件获得这样的
设施~\cite{TimothyHayes2020ARM-HTM,Intel2020TSXdevguide}。
更糟糕的是，一些前沿的调试设施与HTM不兼容~\cite{RobertOCallahan2020DebuggingHTM}。

\subsubsection{弱原子性}
\label{sec:future:Weak Atomicity}

鉴于HTM在可预见的未来可能面临某种大小限制，
HTM有必要与其他机制顺畅地互操作。
如果事务外的读取不会无条件地中止具有冲突写入的事务——
而是简单地为读取提供事务前的值——
那么HTM与风险指针（hazard pointer）和RCU等以读为主的机制之间的互操作性将会得到改善。
这样，风险指针和RCU可以被用来允许HTM处理更大的数据结构并减少冲突概率。

然而，这并不一定简单。
最直接的实现方式需要在每条cache line和总线上增加一个额外的状态，
这是一笔不小的额外开销。
与这笔开销相伴的好处是允许大占用量的读取者，
而不会因持续冲突而饿死更新者。
一种替代方法由Siakavaras等人~\cite{Siakavaras2017CombiningHA}
对二叉搜索树取得了极好的效果，
即使用RCU进行只读遍历，而仅使用HTM进行实际的更新操作。
这种组合的性能比其他事务内存技术高出多达220\pct，
与Howard和Walpole~\cite{PhilHoward2011RCUTMRBTree}
将RCU与STM结合时观察到的加速比相似\@。
在这两种情况下，弱原子性都是在软件中而非硬件中实现的。
尽管如此，看看在硬件和软件中同时实现弱原子性能获得
哪些额外的加速，仍然是一件有趣的事情。

\subsection{结论}
\label{sec:future:Conclusions}

虽然当前的HTM实现在某些情况下确实带来了真正的性能优势，
但它们也有显著的缺陷。
最显著的缺陷似乎是
有限的事务大小、
冲突处理的需求、中止和回滚的需求、
缺乏前进保证、
无法处理不可撤销操作，
以及与加锁之间微妙的语义差异。
关于HTM实现可靠性，也有挥之不去的
担忧~\cite{ChristianJacobi2012MainframeTM,ScottWasson2014HaswellDisableTSX,Intel2020HaswellTSXordering,Intel2021perfTSXordering,MichaelLarabel2021DisableTSX,PowerISA3.1-2020}。

其中一些缺陷可能在未来的实现中得到缓解，
但似乎仍然存在强烈的需求，使HTM能够与许多其他类型的同步机制良好配合，
正如之前所指出的~\cite{McKenney2007PLOSTM,PaulEMcKenney2010OSRGrassGreener}。
虽然已经有一些将HTM与RCU结合使用的
工作~\cite{Siakavaras2017CombiningHA,DimitriosSiakavaras2020RCU-HTM-B+Trees,ChristinaGiannoula2018HTM-RCU-graphcoloring,SeongJaePark2020HTMRCUlock}，
但几乎没有证据表明HTM在与RCU以及其他延迟回收机制的配合方面取得了进展。

简而言之，当前的HTM实现似乎是并行程序员工具箱中受欢迎且有用的补充，
要利用好它们还需要大量有趣且具有挑战性的工作。
然而，它们不能被视为一根魔杖，
用来消除所有并行编程问题。

\QuickQuiz{
	但随着持续的工作，HTM最终实现完整的TM愿景不是必然的吗？
}\QuickQuizAnswer{
	也许吧？

	请回顾\cref{chp:Hardware and its Habits}，尤其是
	对\cref{tab:cpu:CPU 0 View of Synchronization Mechanisms on 8-Socket System With Intel Xeon Platinum 8176 CPUs at 2.10GHz}
	的讨论，它概述了细粒度全局一致的代价。
	还请回顾对\cref{fig:defer:Pre-BSD Routing Table Protected by RCU QSBR}
	的讨论，它展示了避免全局一致需求所带来的可扩展性优势。

	TM主要用于细粒度同步，
	对于每组并发事务，它要求以下全局一致：
	\begin{enumerate*}[(1)]
	\item 事务是有序的，
	或者
	\item 事务之间不冲突。
	\end{enumerate*}
	当然，并发组中可能只有部分事务冲突，
	在这种情况下只有那些冲突的事务需要排序。
	因此，\cref{chp:Hardware and its Habits}中涵盖的内容
	以及对\cref{fig:defer:Pre-BSD Routing Table Protected by RCU QSBR}
	的讨论都直接适用于TM。

	当然，非常聪明的人们确实在继续研究事务内存，
	可以期待他们继续取得进展。
	然而，TM对全局一致的要求与\cref{chp:Hardware and its Habits}中
	讨论的物理定律相冲突，
	这种冲突可能会继续对一般情况下的TM性能和可扩展性施加硬性限制。
}\QuickQuizEnd
